\documentclass[12pt]{article}
\usepackage[utf8]{inputenc}
\usepackage{amsmath}
\usepackage{amssymb}
\usepackage{graphicx}
\usepackage{booktabs}
\usepackage{listings}
\usepackage{xcolor}
\usepackage{hyperref}
\usepackage{geometry}
\geometry{margin=1in}

\lstset{
    language=Python,
    basicstyle=\ttfamily\small,
    keywordstyle=\color{blue}\bfseries,
    commentstyle=\color{green!60!black},
    stringstyle=\color{red},
    numbers=left,
    numberstyle=\tiny\color{gray},
    stepnumber=1,
    numbersep=5pt,
    backgroundcolor=\color{gray!10},
    frame=single,
    breaklines=true,
    breakatwhitespace=false,
    tabsize=4,
    showstringspaces=false
}

\title{Wildfire Costs Calculation Documentation}
\author{}
\date{}

\begin{document}

\maketitle

\section{Introduction}

The \texttt{wildfire\_costs.py} script calculates expected wildfire costs using a probabilistic risk approach. It models ignition rates by terrain and construction type, event severity, and applies a growing annuity formula to account for increasing risk over the project lifetime.

\section{Method Overview}

\subsection{Purpose}

This script calculates wildfire costs by:
\begin{enumerate}
    \item Calculating segment-specific ignition rates by terrain and construction type
    \item Aggregating to total annual event rate (lambda)
    \item Computing Expected Annual Loss (EAL) = lambda × severity
    \item Calculating nominal total cost using growing annuity (accounting for risk growth)
    \item Computing present value using growing annuity formula
\end{enumerate}

\subsection{Calculation Approach}

The wildfire cost calculation follows this structure:

\begin{align}
f_{i,t} &= \text{Base Rate}_{terrain} \times \text{Construction Multiplier} \\
\lambda_{segment} &= \text{Miles}_{terrain} \times f_{i,t} \\
\lambda_{total} &= \sum_{\text{terrains}} \lambda_{segment} \\
\text{EAL} &= \lambda_{total} \times \text{Severity} \\
\text{Nominal Total} &= \text{EAL} \times \frac{(1+g)^N - 1}{g} \quad \text{(if } g \neq 0 \text{)} \\
\text{PV} &= \text{EAL} \times \frac{1 - \left(\frac{1+g}{1+d}\right)^N}{d - g} \quad \text{(if } d \neq g \text{)}
\end{align}

where:
\begin{itemize}
    \item $f_{i,t}$ = effective ignition rate per mile for terrain $i$ and construction type $t$
    \item $\lambda$ = annual event rate (events/year)
    \item $g$ = risk growth rate (annual increase in expected losses)
    \item $d$ = discount rate
    \item $N$ = project lifetime
\end{itemize}

\subsection{Inputs and Outputs}

\textbf{Inputs:}
\begin{itemize}
    \item Project technical details (construction type)
    \item Physical details (terrain distribution)
    \item Wildfire parameters (ignition rates by terrain, severity, growth rate, construction multipliers)
    \item Discount rate source (social, WACC real, or custom)
    \item Project lifetime
\end{itemize}

\textbf{Outputs:}
\begin{itemize}
    \item Total annual event rate (lambda)
    \item Expected Annual Loss (EAL)
    \item Event rates by terrain segment
    \item Nominal total cost (growing annuity)
    \item Present value cost (growing annuity with discounting)
\end{itemize}

\subsection{Integration with Other Scripts}

This script integrates with:
\begin{itemize}
    \item \texttt{yaml\_loaders.py} - Loads project configuration and wildfire parameters (see \texttt{utilities\_documentation.tex})
    \item \texttt{financial\_utils.py} - Financial calculations (see \texttt{utilities\_documentation.tex})
    \item \texttt{csv\_output\_manager.py} - Writes results to CSV (see \texttt{csv\_output\_manager\_documentation.tex})
\end{itemize}

\section{Key Functions}

\subsection{\texttt{get\_discount\_rate(wildfire\_yaml, financing\_yaml)}}

Gets discount rate based on configuration source.

\textbf{Parameters:}
\begin{itemize}
    \item \texttt{wildfire\_yaml} (dict): Loaded wildfire YAML data
    \item \texttt{financing\_yaml} (dict): Loaded financing YAML data
\end{itemize}

\textbf{Returns:}
\begin{itemize}
    \item \texttt{discount\_rate} (float): Discount rate to use
    \item \texttt{source\_description} (str): Description of rate source
\end{itemize}

\textbf{Key Logic:}
\begin{lstlisting}
source = wildfire_yaml["wildfire"]["discount_rate_source"]

if source == "social":
    rate = financing_yaml["financial"]["social_discount_rate"]
elif source == "wacc_real":
    wacc_nominal = financing_yaml["financial"]["wacc_nominal"]
    inflation = financing_yaml["financial"]["inflation_rate"]
    rate = (1 + wacc_nominal) / (1 + inflation) - 1
elif source == "custom":
    rate = wildfire_yaml["wildfire"]["discount_rate_custom"]
\end{lstlisting}

\subsection{\texttt{calculate\_wildfire\_costs(wildfire\_yaml, construction\_type, terrain\_miles, project\_lifetime, discount\_rate)}}

Calculates expected wildfire costs using segment-based ignition rates.

\textbf{Parameters:}
\begin{itemize}
    \item \texttt{wildfire\_yaml} (dict): Loaded wildfire YAML data
    \item \texttt{construction\_type} (str): Construction type (overhead, underground, etc.)
    \item \texttt{terrain\_miles} (dict): Dictionary of terrain type to miles
    \item \texttt{project\_lifetime} (int): Project lifetime in years
    \item \texttt{discount\_rate} (float): Discount rate for PV calculation
\end{itemize}

\textbf{Returns:}
\begin{itemize}
    \item Dictionary containing:
    \begin{itemize}
        \item \texttt{lambda\_total} (float): Total annual event rate
        \item \texttt{EAL} (float): Expected Annual Loss
        \item \texttt{lambda\_by\_terrain} (dict): Event rates by terrain segment
        \item \texttt{severity} (float): Mean loss per event
        \item \texttt{growth\_rate} (float): Risk growth rate
        \item \texttt{nominal\_total} (float): Total nominal cost
        \item \texttt{pv\_cost} (float): Present value cost
    \end{itemize}
\end{itemize}

\textbf{Key Logic:}
\begin{lstlisting}
# Calculate segment-specific ignition rates
for terrain, miles in terrain_miles.items():
    if miles > 0:
        base_rate = ignition_rates_by_terrain.get(terrain, 0.0)
        f_i_t = base_rate * construction_multiplier
        lambda_segment = miles * f_i_t
        lambda_total += lambda_segment

# Calculate EAL
EAL = lambda_total * severity

# Calculate nominal total (growing annuity)
if abs(growth_rate) < 1e-9:
    nominal_total = EAL * project_lifetime
else:
    nominal_total = EAL * ((1 + growth_rate) ** project_lifetime - 1) / growth_rate

# Calculate PV (growing annuity with discounting)
if abs(discount_rate - growth_rate) < 1e-9:
    pv_cost = EAL * project_lifetime
else:
    pv_cost = EAL * ((1 - ((1 + growth_rate) / (1 + discount_rate)) ** project_lifetime) 
                     / (discount_rate - growth_rate))
\end{lstlisting}

\subsection{\texttt{main()}}

Main execution function that orchestrates the wildfire cost calculation and output.

\textbf{Workflow:}
\begin{enumerate}
    \item Loads project technical details and terrain distribution
    \item Loads wildfire parameters
    \item Gets discount rate based on configuration
    \item Calls \texttt{calculate\_wildfire\_costs()} to compute costs
    \item Displays detailed results and writes to CSV
\end{enumerate}

\section{Example}

The following example demonstrates a typical usage of the wildfire cost calculation:

\begin{lstlisting}
# Example: 500 MW Overhead AC line
# Category: "Overhead/AC/500MW/Advanced Aluminum/NA"
# 
# Inputs (from YAML):
#   - Terrain distribution:
#     * Forested: 40 miles, base rate: 0.0001 events/mi/yr
#     * Farmland: 60 miles, base rate: 0.00005 events/mi/yr
#   - Construction multiplier: 1.0 (overhead baseline)
#   - Severity: $50M per event
#   - Risk growth rate: 2% per year
#   - Project lifetime: 50 years
#   - Discount rate: 3% (social)
#
# Calculation:
#   Forested segment: 40 * 0.0001 * 1.0 = 0.004 events/year
#   Farmland segment: 60 * 0.00005 * 1.0 = 0.003 events/year
#   Total lambda = 0.007 events/year
#   EAL = 0.007 * $50M = $350,000/year
#   Nominal total = $350k * ((1.02^50 - 1) / 0.02) = $28.5M
#   PV = $350k * ((1 - (1.02/1.03)^50) / (0.03 - 0.02)) = $12.3M
#
# Output:
#   - Total annual events: 0.007 events/year
#   - EAL: $350,000/year
#   - Nominal total: $28.5M
#   - Present value: $12.3M
\end{lstlisting}

\textbf{Integration Example:}

\begin{lstlisting}
# When called from ctcc.py:
csv_manager = CTCCOutputManager()
# Results dict already has all needed values from calculate_wildfire_costs()
csv_manager.add_wildfire_costs(results)
csv_manager.write_batch_summary()
\end{lstlisting}

\section{Key Implementation Details}

\subsection{Probabilistic Risk Approach}

The script uses a probabilistic approach where wildfire events are modeled as Poisson processes. The annual event rate (lambda) represents the expected number of events per year, and EAL represents the expected annual loss.

\subsection{Construction Type Multipliers}

Different construction types have different ignition rates. Overhead lines have the highest risk (multiplier = 1.0), while underground and subsea have much lower multipliers (e.g., 0.01 for underground).

\subsection{Growing Annuity Formula}

The script accounts for increasing wildfire risk over time using a growing annuity formula. This reflects the reality that climate change and other factors may increase wildfire frequency and severity.

\subsection{Edge Cases}

The script handles edge cases:
\begin{itemize}
    \item If growth rate = 0: Uses simple annuity formula
    \item If discount rate = growth rate: Uses simplified formula (PV = EAL × N)
\end{itemize}

\subsection{AFUDC Treatment}

Wildfire costs are probabilistic future losses, not capital costs. Therefore, they are not AFUDC-eligible. AFUDC applies only to construction work in progress.

\subsection{Warning for High Growth Rates}

If the risk growth rate exceeds the discount rate, the script issues a warning because the present value grows over time, indicating that risk is escalating faster than discounting.

\end{document}

